% answers_short.tex —— 2011 年·重庆卷·文科 · 简短答案速查
% 由 tools/make_short_answers.py 自 answers.tex 的 \gksolhead/\gkanswer 提取，
% 并与 2011/2011_数学真题_重庆文_答案要点.md 逐题比对后生成，勿手改（改源 MD 后重跑）。

\begin{gknotice}
\textbf{使用说明}：本节只给最终答案便于快速核对，完整推导见解析版。答案为编写者独立推导，\textbf{非官方标准答案}。
\end{gknotice}

\gkpart{选择题}
\begin{center}\begin{tabular}{ccccccccccc}
\toprule
题号 & 1 & 2 & 3 & 4 & 5 & 6 & 7 & 8 & 9 & 10 \\
\midrule
答案 & D & A & A & C & D & B & C & D & B & A \\
\bottomrule
\end{tabular}\end{center}

\gkpart{填空题}
\begin{itemize}[leftmargin=2.6em]
  \item[\textbf{11.}] $240$。
  \item[\textbf{12.}] $\dfrac43$。
  \item[\textbf{13.}] $y=2x$。
  \item[\textbf{14.}] $\dfrac7{30}$。
  \item[\textbf{15.}] $2-\log_23$。
\end{itemize}

\gkpart{解答题}
\begin{itemize}[leftmargin=2.6em]
  \item[\textbf{16.}] (1) $a_n=2^n$；(2) $S_n=2^{n+1}+n^2-2$
  \item[\textbf{17.}] (1) $\dfrac{16}{81}$；(2) $\dfrac49$
  \item[\textbf{18.}] (1) $\pi$；(2) $\dfrac{3\sqrt3}2$
  \item[\textbf{19.}] (1) $a=3$，$b=-12$；(2) 极大值 $f(-2)=21$，极小值 $f(1)=-6$
  \item[\textbf{20.}] (1) $\dfrac{\sqrt5}8$；(2) $\dfrac{2\sqrt{15}}7$
  \item[\textbf{21.}] (1) $\dfrac{x^2}4+\dfrac{y^2}2=1$；(2) 存在，$F(\sqrt{10},0)$，比值为 $\dfrac{\sqrt2}2$
\end{itemize}
